% Источники: exam/materials/моделирование.pdf, раздел 8.
\section{Системы, меняющие своё поведение во времени и гибридные сложные системы.}

Одной из черт сложного поведения является наличие у системы нескольких
качественно различных, последовательно сменяющих друг друга во времени,
поведений. Инженеры обычно называют их \textbf{режимами функционирования}.

Такой тип сложного поведения можно реализовать, если описать всю
совокупность допустимых, простых, в некотором смысле, частных поведений и
указать правила переключения с одного поведения на другое. Организованная
таким образом сложная динамическая система в каждый конкретный момент
времени ведет себя как некоторая простая динамическая система.

Каждое конкретное поведение можно отождествить со значением некоторой
дискретной переменной, а мгновенные переключения текущего поведения -- с
дискретными событиями. Для передачи информации о дискретных событиях в
другие блоки используют специальные переменные -- сигналы.

Набор дискретных состояний вместе с условиями переходов из одного состояния
в другое образует обычное дискретное поведение. В моменты переходов могут
происходить мгновенные скачкообразные изменения значений переменных.
Поскольку в каждом из дискретных состояний элементарный блок ведет себя как
некоторая непрерывная система, то поведение блока в целом является
\textbf{непрерывно-дискретным} или \textbf{гибридным}.

Пример из источника: бассейн с двумя трубами обладает сложным поведением,
поскольку его поведение при переполнении качественно отличается от поведения
при нормальном уровне и при опустошении. Такая система может быть
представлена как совокупность трех простых динамических систем: нормальный
уровень, переполненный бассейн и пустой бассейн.

Объединяя структурно-сложные системы и гибридные, получаем новый тип
сложных систем -- \textbf{структурно-сложные гибридные системы}. Их главная
черта -- параллельное функционирование нескольких гибридных систем
иерархической структуры.

\clearpage
